%=============================================================
% POSTER ACADÉMICO — DIAPCE
% Distribución: Título · Autores · Introducción · Objetivos ·
%               Marco Teórico · Metodología · Resultados ·
%               Conclusiones · Referencias
% Tamaño: A1 horizontal (841 × 594 mm)
%=============================================================
\documentclass[10pt]{article}

% === GEOMETRÍA ===
\usepackage[
  a1paper,
  landscape,
  top=1.2cm, bottom=1.2cm,
  left=1.2cm, right=1.2cm
]{geometry}

% === PAQUETES ESTÁNDAR ===
\usepackage[utf8]{inputenc}
\usepackage[T1]{fontenc}
\usepackage[spanish]{babel}
\usepackage{graphicx}
\usepackage{amsmath}
\usepackage{xcolor}
\usepackage[most]{tcolorbox}
\usepackage{enumitem}
\usepackage{url}
\usepackage{helvet}          % Fuente sans-serif
\renewcommand{\familydefault}{\sfdefault}
\usepackage{microtype}

\pagestyle{empty}
\setlength{\parindent}{0pt}

%=============================================================
% PALETA DE COLORES (inspirada en la plantilla del poster)
%=============================================================
\definecolor{pinkprimary}{RGB}{232, 58, 112}    % Rosa principal (header y títulos)
\definecolor{pinksecondary}{RGB}{248, 187, 208} % Rosa claro (subtítulos header)
\definecolor{boxbg}{RGB}{250, 250, 250}         % Fondo de cajas
\definecolor{accentteal}{HTML}{2E86AB}           % Azul-verde para Resultados
\definecolor{darktext}{RGB}{30, 30, 50}

%=============================================================
% ESTILOS DE CAJAS (tcolorbox)
%=============================================================
\tcbset{
  posterbox/.style={
    enhanced,
    colback=boxbg,
    colframe=pinkprimary,
    coltitle=white,
    fonttitle=\bfseries\normalsize,
    attach boxed title to top left={xshift=6mm, yshift=-3.2mm},
    boxed title style={
      colback=pinkprimary,
      colframe=pinkprimary,
      arc=3pt
    },
    arc=4pt,
    boxrule=1.2pt,
    top=4mm, left=3mm, right=3mm, bottom=2mm,
    before skip=0mm, after skip=3mm
  },
  tealbox/.style={
    posterbox,
    colframe=accentteal,
    boxed title style={colback=accentteal, colframe=accentteal, arc=3pt}
  },
  subbox/.style={
    posterbox,
    fonttitle=\bfseries\small,
    top=3mm, bottom=2mm,
    before skip=0mm, after skip=0mm
  }
}

% Conector de flujo entre sub-cajas
\newcommand{\flowconn}[1]{%
  \vspace{0.5mm}%
  \begin{center}%
    \textcolor{#1}{\rule{2cm}{0.8pt}\enspace$\downarrow$\enspace\rule{2cm}{0.8pt}}%
  \end{center}%
  \vspace{0.5mm}%
}

%=============================================================
\begin{document}

%-------------------------------------------------------------
%  ENCABEZADO: TÍTULO + AUTORES Y AFILIACIÓN
%-------------------------------------------------------------
\begin{tcolorbox}[
  enhanced,
  colback=pinkprimary,
  colframe=pinkprimary,
  arc=6pt, boxrule=0pt,
  top=10mm, bottom=8mm, left=14mm, right=14mm
]
  \centering
  {\fontsize{32}{38}\selectfont\bfseries\color{white}
    DIAPCE: Diseño y Análisis de Proporciones de Concreto Estructural%
  }\\[5mm]
  {\fontsize{15}{19}\selectfont\color{pinksecondary}
    Software de implementación de un modelo matemático para el análisis y
    optimización de mezclas de concreto estructural — Meta, Colombia%
  }\\[7mm]
  {\fontsize{13}{17}\selectfont\bfseries\color{white}
    Juan Diego Pedroza Sierra%
  }\\[3mm]
  {\fontsize{11}{15}\selectfont\color{pinksecondary}
    Universidad Cooperativa de Colombia\enspace·\enspace
    Facultad de Ingeniería\enspace·\enspace
    Programa de Ingeniería de Sistemas\enspace·\enspace
    Villavicencio, Meta, Colombia\enspace·\enspace 2023%
  }
\end{tcolorbox}

\vspace{5mm}

%-------------------------------------------------------------
%  CUERPO: 4 COLUMNAS
%  Col 1 → Introducción + Objetivos
%  Col 2 → Marco Teórico
%  Col 3 → Metodología
%  Col 4 → Resultados + Conclusiones + Referencias
%-------------------------------------------------------------
\noindent
%==============================
% COLUMNA 1
%==============================
\begin{minipage}[t]{0.2425\textwidth}

  \begin{tcolorbox}[posterbox, title=Introducción]
    \small
    El departamento del Meta (Colombia) presenta condiciones ambientales
    agresivas —altas temperaturas y humedad elevada— que aceleran la degradación
    del concreto y elevan los costos de mantenimiento.

    En los \textbf{laboratorios de materiales y semilleros de investigación de
    la Universidad Cooperativa de Colombia (UCC)} persistía una brecha de
    herramientas digitales locales para apoyar la dosificación, la selección de
    aditivos y la trazabilidad de ensayos académicos de forma reproducible.

    \textbf{DIAPCE} es una aplicación multiplataforma en Flutter
    \textbf{dirigida exclusivamente a los estudiantes del semillero y laboratorio
    de la UCC}; no es una herramienta de uso general ni industrial. Opera sin
    conexión (offline-first) mediante SQLite, y complementa —sin sustituir— el
    criterio del docente y el marco normativo, aportando evidencia sistematizada
    de carácter formativo.
  \end{tcolorbox}

  \begin{tcolorbox}[posterbox, title=Objetivos]
    \small
    \textbf{General}\\
    Desarrollar un software que implemente un modelo matemático para el ajuste
    de proporciones de materiales cementantes en mezclas de concreto estructural.

    \vspace{2pt}
    \textbf{Específicos}
    \begin{enumerate}[leftmargin=14pt, topsep=1pt, itemsep=1pt]
      \item Implementar una base de datos para el almacenamiento de
            características físicas, mecánicas y ambientales de mezclas de
            concreto estructural.
      \item Desarrollar el software con un modelo de decisión en cascada para
            dosificación y clasificación de aditivos (P0–PP3), proyectando
            resistencias a 7, 14 y 28 días según ACI\,211.1.
      \item Realizar pruebas y análisis de desempeño del programa aplicando
            criterios de calidad y métricas de clasificación y error de
            proyección.
    \end{enumerate}
  \end{tcolorbox}

\end{minipage}
\hfill
%==============================
% COLUMNA 2
%==============================
\begin{minipage}[t]{0.2425\textwidth}

  \begin{tcolorbox}[subbox, title=Marco Teórico — Ley de Abrams / ACI\,211.1]
    \small
    Resistencia $f'_c$ inversa a la relación a/c (ley de Abrams):
    \[ f'_c = K\cdot\!\ \left(\tfrac{w}{c}\right)^{-n} \]
    ACI\,211.1 traduce esto en proporciones prácticas.
    En DIAPCE la a/c se restringe al dataset experimental.
  \end{tcolorbox}

  \flowconn{pinkprimary}

  \begin{tcolorbox}[subbox, title=Marco Teórico — Condiciones Ambientales]
    \small
    $T$ y HR condicionan la hidratación. Flujo de decisión DIAPCE:
    \[ T^{\circ}\;{\rightarrow}\;HR\;{\rightarrow}\;a/c\;{\rightarrow}\;\text{aditivo} \]
    Contexto: Meta, $T$ 25–35\,°C, régimen higrométrico variable.
  \end{tcolorbox}

  \flowconn{pinkprimary}

  \begin{tcolorbox}[subbox, title=Marco Teórico — Aditivos P0–PP3 (ASTM\,C494)]
    \small
    \textbf{P0} control\enspace·\enspace
    \textbf{P1/P2/P3} impermeabilizante \textit{(Euco Vandex AM\,10I)} 2/3/4\,\% s.m.c.\\
    \textbf{PP1/PP2/PP3} plastificante \textit{(Sika\textsuperscript{\textregistered} Plastiment\textsuperscript{\textregistered} AP)} 0{,}2\,/\,0{,}4\,/\,0{,}6\,\% s.m.c.\\
    Solo se muestran los códigos válidos para la condición seleccionada.
  \end{tcolorbox}

  \flowconn{pinkprimary}

  \begin{tcolorbox}[subbox, title=Marco Teórico — Proyección y Modelado]
    \small
    Resistencias 7/14/28\,días (ASTM\,C39):
    \[ \hat{f}_{t}=\operatorname{avg}\{f_{t,i}:\text{cond.}(i)\in\text{sel.}\} \]
    Modelo empírico: transparente, robusto, compatible con la cascada.
    Integración futura de ML sobre las mismas vistas SQL.
  \end{tcolorbox}

\end{minipage}
\hfill
%==============================
% COLUMNA 3
%==============================
\begin{minipage}[t]{0.2425\textwidth}

  \begin{tcolorbox}[subbox, title=Metodología — Enfoque y Dataset]
    \small
    \textbf{FDD} (Feature-Driven Development), ágil, investigación mixta.\\
    \textbf{1\,874 registros} regionales del Meta\,·\,24–57\,MPa\,·\,datos desde CSV.
  \end{tcolorbox}

  \flowconn{pinkprimary}

  \begin{tcolorbox}[subbox, title=Metodología — Flujo de Decisión]
    \small
      \begin{enumerate}[leftmargin=14pt, topsep=0pt, itemsep=0pt]
      \item Resistencia objetivo (24–57\,MPa).
      \item Temperatura (°C). \quad 3.\,Humedad (\%).
      \item[4.] Filtra a/c válidos.
      \item[5.] Aditivo: \textbf{P0}–\textbf{P3} / \textbf{PP1}–\textbf{PP3}.
      \item[6.] Calcula $\hat{f}_7$, $\hat{f}_{14}$, $\hat{f}_{28}$.
      \item[7.] Guarda en SQLite (trazabilidad).
    \end{enumerate}
  \end{tcolorbox}

  \flowconn{pinkprimary}

  \begin{tcolorbox}[subbox, title=Metodología — Base de Datos SQLite]
    \small
    4 tablas: \texttt{TiposAditivo}\,·\,\texttt{Productos}\,·\,\texttt{Aditivos}\,·\,\texttt{ResultadosConcreto}.\\
    FK + borrado en cascada\,·\,índice compuesto (temp, humid, a/c, aditivo, edad).
  \end{tcolorbox}

  \flowconn{pinkprimary}

  \begin{tcolorbox}[subbox, title=Metodología — Vistas de la App]
    \small
    \textbf{Login/Registro}\,—\,auth local\enspace·\enspace
    \textbf{Hall}\,—\,panel de proyectos.\\
    \textbf{Crear}\,—\,selección en cascada con validación de rangos.\\
    \textbf{Detalle}\,—\,predicciones $\hat{f}_{7/14/28}$ + pastel + línea (\texttt{fl\_chart}).
  \end{tcolorbox}

\end{minipage}
\hfill
%==============================
% COLUMNA 4
%==============================
\begin{minipage}[t]{0.2425\textwidth}

  \begin{tcolorbox}[tealbox, title=Resultados]
    \small
    \begin{itemize}[leftmargin=12pt, topsep=1pt, itemsep=1pt]
      \item \textbf{App funcional offline-first:} Login, registro, panel de
            proyectos, creación y consulta, operativos sin conexión.
      \item \textbf{Predicciones en tiempo real:} Promedios $\hat{f}_7$,
            $\hat{f}_{14}$ y $\hat{f}_{28}$ calculados desde 1\,874 registros
            locales en rango 24–57\,MPa.
      \item \textbf{Clasificación guiada P0–PP3:} Impermeabilizantes
            \textbf{P1}–\textbf{P3} (Euco Vandex AM\,10I, 2/3/4\,\%) y
            plastificantes \textbf{PP1}–\textbf{PP3} (Sika Plastiment AP,
            0{,}2/0{,}4/0{,}6\,\%) filtrados por condición ambiental;
            evita selecciones sin respaldo experimental.
      \item \textbf{Visualizaciones:} Gráfico de pastel (composición de mezcla)
            y gráfico de línea (evolución de resistencia a 7/14/28 días).
      \item \textbf{Trazabilidad:} Cada proyecto registra parámetros,
            código de aditivo, predicciones y fecha en SQLite.
      \item \textbf{Multiplataforma:} Android, iOS y Web desde un único
            código base Flutter/Dart.
    \end{itemize}
  \end{tcolorbox}

  \begin{tcolorbox}[posterbox, title=Conclusiones]
    \small
    \begin{itemize}[leftmargin=12pt, topsep=1pt, itemsep=1pt]
      \item \textbf{Herramienta formativa exclusiva:} DIAPCE está dirigida a
            \textbf{estudiantes del semillero y laboratorio de la UCC, Meta};
            su alcance es institucional y académico, no industrial.
      \item \textbf{Trazabilidad académica:} Registro reproducible de insumos,
            resultados y versión del modelo, favoreciendo evaluación docente.
      \item \textbf{Validez normativa:} Clasificación P0–PP3 y proyecciones
            7/14/28 días coherentes con ACI\,211.1 en contexto formativo.
      \item \textbf{Adopción sin fricciones:} Diseño offline-first operativo
            en entornos con conectividad limitada y equipos heterogéneos.
      \item \textbf{Trabajo futuro:} Ampliar el dataset regional, explorar
            modelos supervisados (ML) y agregar módulo de sincronización
            cliente-servidor.
    \end{itemize}
  \end{tcolorbox}

  \begin{tcolorbox}[posterbox, title=Referencias Bibliogr\'{a}ficas]
    \scriptsize
    \begin{itemize}[leftmargin=10pt, topsep=0pt, itemsep=1pt]
      \item ACI\,211.1-16. \textit{Standard Practice for Selecting Proportions
            for Normal Concrete.} American Concrete Institute, 2016.
      \item ASTM C39/C39M. \textit{Compressive Strength of Cylindrical Concrete
            Specimens.} ASTM International.
      \item ASTM C494/C494M-22. \textit{Chemical Admixtures for Concrete.}
            ASTM International, 2022.
      \item Palmer, S.\,R. \& Felsing, J.\,M. (2002). \textit{A Practical Guide
            to Feature-Driven Development.} Prentice Hall.
      \item Flutter.dev. \textit{Flutter — Build apps for any screen.}
            \url{https://flutter.dev}
      \item SQLite Consortium. \textit{SQLite Documentation.}
            \url{https://www.sqlite.org}
      \item García, P. \& Torres, R. (2021). Integración de materiales
            cementantes alternativos. \textit{J. Sustainable Construction}, 8(3).
      \item Departamento Nacional de Planeación. (2018).
            \textit{Política Nacional de Edificaciones Sostenibles.} DNP.
    \end{itemize}
  \end{tcolorbox}

\end{minipage}

\end{document}
